\chapter{Testing and validation}
\label{chap:testing}

\section{The V-model and the four test levels}
For the testing of RAM-USB the V-model was adopted, an extension of the waterfall model in which
each specification phase is paired with a corresponding verification phase, written hand in hand with
the specification.

This approach was adopted because RAM-USB was deliberately treated as a
safety-critical system, in which a defect in production has serious consequences for the users of the system.
It was therefore essential that an error in the specification be detected well before implementation and
that testing be applied at several specification levels.

\autoref{tab:livelli-test} shows the correspondence between the requirement category and the test
level that verifies it, along with the technique used at that level.

\begin{table}[H]
\centering
\caption{Verification level associated with each specification level.}
\label{tab:livelli-test}
\begin{tabular}{@{}l p{\dimexpr0.55\linewidth\relax}@{}}
\hline
\textbf{Specification level} & \textbf{Verification level} \\
\hline
Component requirement & Unit testing with Go's \texttt{testing} package \\
Use case & Integration testing, bottom-up with drivers and stubs \\
Non-functional requirement & System testing, on the whole stack \\
User requirement & Acceptance checklist, walked through by hand on the corresponding use case \\
\hline
\end{tabular}
\end{table}

\section{Bottom-up integration strategy}

Integration starts from the innermost component and proceeds outward, in the order Database-Vault,
Storage-Service, Security-Switch, Network-Manager, Entry-Hub.

The choice to start from Database-Vault is dictated by the fact that it is at the same time the component
with the fewest outgoing dependencies and the one on which the largest number of functional requirements is concentrated.
Building it first, with drivers standing in for the components that invoke it and a stub standing in for
Storage-Service, the only component it invokes, ensures that the most sensitive component is also the one verified the longest.


\section{TDD: entry and exit criteria}

The entry criterion is that a test exists for every requirement, and that any failure of that test
occurs before the implementation of the requirement under test.

An example: requirement DV-F-15 requires that a nonexistent email and a wrong password produce the
same login outcome, indistinguishable from the outside.

The test that verifies it, \texttt{TestLogin\_NonexistentEmailAndWrongPassword\_AreIndistinguishable}
\cite{ramusb-code-login-test-dv-f-15}, calls \texttt{Login} with the two different credentials and requires
that the outcomes match exactly.

A test that has never been seen to fail, perhaps written after the code and passing on the first try, offers no
guarantee that it verifies anything: it might compare two values that are equal by construction, or not call the
function it is supposed to test at all.
For this reason, the test is written before the body of the \texttt{Login} function actually exists, and
it fails, providing proof that the test can recognize a wrong outcome.

Only after the test has failed does it make sense to write the complete body of \texttt{Login}, and only then does the
transition from failed test to passed test become useful information.

Writing tests first also helps to disambiguate and clarify requirements. Indeed, if a test cannot be written for a
requirement, the requirement is unclear or ambiguous and
consequently must be reworded more precisely.

There are three exit criteria, and they depend on the code under test. The first concerns only the validation
library, which requires more thorough tests because it is considered particularly sensitive.

\begin{itemize}
    \item in the validation library shared by Entry-Hub, Security-Switch, and Database-Vault
    (\autoref{sec:ciclo-richiesta}), every branch of every \texttt{if} or \texttt{switch} that decides whether to
    accept or reject a field must be traversed by at least one test case;
    \item no test that has already passed on the already integrated chain of components may fail;
    \item a change reruns not only the test written specifically for it but the entire test suite.
\end{itemize}

\section{Coverage and gaps in the test plan}

\texttt{go test -cover} measures, for each package of the codebase, the percentage of statements executed
at least once: \texttt{pkg/validation}, for example, reaches 98\%, because that is where the application logic lives,
while the \texttt{cmd} packages, which start each service, stay at 10--20\%, because verifying them is
the job of integration and system testing.

The exit criterion mentioned above for \texttt{pkg/validation} is a stricter measure that adds to
this one: it is not enough to execute almost all statements, every single branch of every \texttt{if} or
\texttt{switch} must have been taken by at least one test.

The test plan, however, is not free of limitations:
No test is planned for either protecting backups from tampering or rotating the
master key, currently out of scope, but essential in a project of this kind.
Load and stress testing, moreover, are planned but not yet designed. There are therefore no
reliable benchmarks of the system's behavior.
Every end-to-end test on the host involved a single user with a single small file in the backup, so neither
concurrency nor mutual isolation between two users has been put to the test.

\section{Automated pre-commit checks}

In addition to the entry and exit criteria, an automated verification chain runs before every
commit \cite{ramusb-contributing}. \texttt{golangci-lint} \cite{golangci-lint} is a runner that aggregates
multiple linters in a single pass; the RAM-USB configuration enables fifteen of them in total.

Among these, \texttt{gosec} \cite{gosec} looks for known security patterns in the code, and has been present since the
early stages of implementation.

A second extension introduced four linters targeting unreleased resources and mishandled errors:

\begin{itemize}
    \item \texttt{bodyclose} \cite{bodyclose} flags an HTTP response whose body is never
    closed, which therefore leads to a resource leak;
    \item \texttt{sqlclosecheck} \cite{sqlclosecheck} does the same for SQL rows and statements left
    open;
    \item \texttt{contextcheck} \cite{contextcheck} detects a function that does not propagate the
    \texttt{context.Context} received from the caller;
    \item \texttt{errorlint} \cite{errorlint} checks that errors are compared in the
    correct way.
\end{itemize}

These are complemented by Go's race-condition detector, invoked with \texttt{go test -race}.

The chain is completed by \texttt{govulncheck} \cite{govulncheck}, which compares the dependencies against a database
of known vulnerabilities but reports a problem only if the code actually reaches the vulnerable function,
and by two scanners designed to overlap with each other:
\texttt{gitleaks} \cite{gitleaks}, which looks for secrets already committed to the repository history, and
\texttt{trivy} \cite{trivy}, invoked in filesystem scan mode (\texttt{trivy fs .}), which
checks both dependencies and secrets independently of the two previous tools.

The redundancy of the last two is the application of the defense-in-depth concept to testing.

\section{Results: completed verifications}

The most significant campaign was run on real hardware, a Proxmox host with the components
deployed according to the placement described in \autoref{chap:deployment}.

The four use cases passed: registration and login succeeded, the backup completed and the
restore returned a file identical to the original.

\autoref{tab:verifiche} shows the properties checked on the running system.

\begin{table}[H]
\centering
\caption{Properties verified on the running system.}
\label{tab:verifiche}
\begin{tabular}{@{}p{\dimexpr0.34\linewidth\relax} p{\dimexpr0.56\linewidth\relax}@{}}
\hline
\textbf{Property} & \textbf{Observed result} \\
\hline
Hashing parameters & The string stored in the record shows \texttt{argon2id}, version 19,
memory 47104 KiB, two passes, parallelism one \\
Filesystem isolation & Confinement root owned by root with permissions 711,
user subdirectory 700 \\
Identity on the mesh & Every node is owned by the synthetic user
\texttt{tagged-devices}, never by the account created for the person \\
Grant duration & Expiry recorded at exactly twelve hours from the moment of login \\
Reachability boundaries & No node lists among its peers a component that the
policy forbids it from contacting \\
Restart survival & A service that is stopped and restarted comes back into service reusing
the identity on disk, without presenting the already consumed single-use token again \\
\hline
\end{tabular}
\end{table}

During the execution of the complete test plan, several defects emerged that had not been
caught by the earlier unit tests \cite{ramusb-known-issues}.

Two concerned the client and prevented backup entirely. The first was the SSH port: the client
connected to the default one, while the server listens on another. The second was the absence of
any host key policy: \texttt{ssh}, invoked by Restic without a terminal, had no
\texttt{known\_hosts} to check and no way to ask for confirmation, and refused the connection instead of
asking, so the first backup on every new machine always failed. A third concerned the configuration,
because the client has a single address for a service that exposes two on different listeners. A
fourth was a usage string documenting a syntax that the argument parsing cannot accept.

A fifth was more serious: the mTLS identity was not preserved from one startup to the next, so
no service survived the restart of its own container, because it retried the bootstrap with a
single-use token already spent.

Other small logic and implementation errors, too numerous to be reported in this thesis,
were detected and fixed thanks to the tests.
